%% GEP — KDD-style (ACM sigconf) two-column submission draft.
%% Compiles on Overleaf/TeX Live with the `acmart` class (standard).
%% DRAFT: cells marked \tbd{} are NOT yet measured — fill from your own runs.
%% Verify all \cite references against the actual papers before submitting.
\documentclass[sigconf,nonacm]{acmart}

\settopmatter{printacmref=false}
\setcopyright{none}
\pagestyle{plain}

\usepackage{booktabs}
\usepackage{xcolor}
\newcommand{\tbd}{\textcolor{red}{\textbf{TBD}}}

\begin{document}

\title{Grounded Execution Precision for Small Financial Language Models}
\subtitle{Relocating the bottleneck from evidence access to auditable extraction and execution}

\author{Anonymous Author(s)}
\affiliation{\institution{Anonymous Institution}\country{}}

\begin{abstract}
Financial question answering over regulatory filings is widely treated as an evidence-access
problem, tackled with retrieval and long context. We argue and empirically show that for small
and mid-sized open language models the dominant bottleneck is instead \emph{decision extraction
and computation}: even when the exact gold evidence is supplied, small models rarely produce the
correct numeric answer. We introduce \emph{Grounded Execution Precision} (\textsc{GEP}), a
decomposed procedure in which the
model extracts typed decision variables and a formula from the raw filing (never the gold
answer), a deterministic executor computes the result, and a grounding/self-verification signal
gates a selective answer-or-abstain decision. On the FinanceBench open subset, at fixed model
size and under an LLM-as-judge scorer, the procedure raises accuracy on computational questions
from $\sim$20\% (raw-evidence prompting) to $\sim$62\% when it commits (i.e., 62\% at 43\%
coverage), a $\sim$3$\times$ improvement driven purely by the decision procedure. We further show
that the residual error is dominated by generation instability and question-type mismatch, not
reasoning, and outline the distillation study that transfers this behavior into a small student.
\end{abstract}

\keywords{financial NLP, tool use, calibration, selective prediction, evidence grounding}

\maketitle

\section{Introduction}
Large language models are increasingly applied to financial decision support over primary
documents such as 10-K/10-Q filings. The prevailing framing is one of \emph{access}: give the
model the right passages (via retrieval or long context) and it will answer. We provide evidence
that, for small and mid-sized open models, this framing is misleading. Even under an
\emph{oracle} setting---the gold evidence supplied directly---such models answer only $\sim$20\%
of computational questions correctly, while answering near-perfectly when the final value is
handed to them verbatim. The gap is therefore one of \emph{extracting decision variables and
computing over them}, not of retrieval.

\textbf{Hypothesis.} \emph{For small/mid open LLMs on grounded financial QA, the bottleneck is
decision extraction and computation, not evidence access. Decomposing the task---extract typed
variables and a formula, compute with a deterministic tool, and abstain unless the decision is
grounded and executable---substantially improves accuracy and calibration on computational
questions, where raw-evidence prompting does not.}

\textbf{Contributions.}
\begin{itemize}
\item A decomposed, no-leakage procedure (\textsc{GEP}) that turns financial QA into
typed-variable extraction plus deterministic computation, with a grounding/verification
confidence for selective prediction (\S\ref{sec:method}).
\item An apples-to-apples measurement showing a $\sim$3$\times$ accuracy gain at fixed model size
on computational FinanceBench questions, triangulated across three scorers including
LLM-as-judge (\S\ref{sec:results}).
\item A quantified failure taxonomy separating structural task-type limits and infrastructure
instability from reasoning error, and a distillation study design that transfers the behavior to
a small student (\S\ref{sec:results}, \S\ref{sec:limitations}).
\end{itemize}

\section{Related Work}
\textbf{Financial QA benchmarks.} FinanceBench~\cite{financebench} evaluates QA over real filings
and reports that models struggle, especially without oracle context. \textbf{Tool / program-of-
thought reasoning.} Delegating arithmetic to executable programs improves numeric
reasoning~\cite{pot}. \textbf{Self-consistency and calibration.} Sampling agreement and
verbalized confidence estimate correctness~\cite{selfconsistency,calibration}. \textbf{Selective
prediction} trades coverage for accuracy via abstention~\cite{selective}. \textbf{Knowledge
conflict} studies whether models follow provided evidence over parametric
memory~\cite{knowledgeconflict}. \textbf{Rationale distillation} transfers reasoning to smaller
models~\cite{distill}. \textsc{GEP} combines these into a finance-specific, verifiable
decision procedure and studies its distillation, rather than any single ingredient in isolation.

\section{Method}\label{sec:method}
Given a question $x$ and raw evidence $E$ from a filing, \textsc{GEP} factorizes the answer
as extraction followed by deterministic computation:
\[
P(y \mid x, E) \;=\; \sum_{z} P_\theta(z \mid x, E)\, \mathbf{1}[y = f(z)],
\]
where $z=(\mathcal{V},\phi)$ is a program---typed decision variables $\mathcal{V}$ and a formula
$\phi$---and $f$ is a deterministic executor. The extractor sees only $x$ and $E$ (never the gold
answer), emitting $z$ as JSON. The executor evaluates $\phi$ over $\mathcal{V}$ (a sandboxed
arithmetic evaluator); a single-variable program with no formula is treated as a lookup.

\textbf{Grounding-calibrated confidence.} We define a confidence
$c = g(\mathcal{V},E)\cdot\mathbf{1}[\text{executed}]$, where $g$ is the fraction of input values
found verbatim in $E$. A fully grounded, executable program is far more likely correct; a
hallucinated or non-executing one is not. Optionally a stronger self-verification pass rates the
computed answer. The system \emph{abstains} when it cannot produce a grounded executable program;
answering only high-$c$ cases yields selective accuracy at chosen coverage.

\textbf{No-leakage.} Programs are built from raw evidence and question metadata only; the gold
answer never enters the prompt. Unit-tolerant scoring treats scale-equivalent numbers as equal
(e.g., \$400M $\equiv$ \$400{,}000{,}000).

\section{Experimental Setup}\label{sec:setup}
\textbf{Data.} FinanceBench open subset (150 examples). \textbf{Models.} To test generality
across model families (not a single-model artifact) we evaluate Qwen2.5-Instruct
(0.5B/1.5B/3B/7B), DeepSeek (\texttt{deepseek-llm-7b-chat}), NVIDIA Nemotron
(\texttt{Nemotron-Mini-4B-Instruct}), and Phi-3.5-mini, with Gemma-2 and Llama-3.1 as gated
extras---all in fp16 on Apple Silicon (no quantization).
\textbf{Baselines.} question-only and raw gold-evidence prompting at matched model size.
\textbf{Scorers.} strict weak-match, unit-tolerant numeric match, and LLM-as-judge (a stronger,
different-family judge); we report judge-vs-human agreement on a labeled subset. \textbf{Decode.}
greedy (repetition penalties corrupt JSON extraction). Pilot results use Qwen2.5-7B, $n{=}30$.

\section{Results}\label{sec:results}

\begin{table}[t]
\caption{Main result (Qwen2.5-7B, FinanceBench, $n{=}30$ pilot). Same model, same scorer.
``@commit'' is accuracy on answered items (abstention = selective prediction).}
\label{tab:main}
\begin{tabular}{lcccc}
\toprule
Method & Cov. & Acc@full & Acc@commit & Exec. \\
\midrule
Raw gold-evidence (baseline) & 1.00 & 0.20 & 0.20 & --- \\
\textbf{\textsc{GEP} (ours)} & 0.43 & 0.27 & \textbf{0.62} & 0.43 \\
\bottomrule
\end{tabular}
\end{table}

\begin{table}[t]
\caption{Selective accuracy vs.\ coverage (self-verification ranking).}
\label{tab:selective}
\begin{tabular}{lcccc}
\toprule
Coverage & 20\% & 30\% & 50\% & 100\% \\
\midrule
Accuracy & \textbf{0.83} & 0.56 & 0.40 & 0.20 \\
\bottomrule
\end{tabular}
\end{table}

\begin{table}[t]
\caption{Accuracy by question type (unit-tolerant, $n{=}30$).}
\label{tab:tasktype}
\begin{tabular}{lccc}
\toprule
Task type & Acc. & Exec. & $n$ \\
\midrule
guidance\_delta & 1.00 & 1.00 & 1 \\
line\_item\_lookup & 0.60 & 0.60 & 5 \\
ratio\_calculation & 0.27 & 0.64 & 11 \\
period\_comparison & 0.17 & 0.33 & 6 \\
cash\_flow\_category\_selection & 0.00 & 0.00 & 1 \\
generic\_financial\_qa & 0.00 & 0.00 & 6 \\
\bottomrule
\end{tabular}
\end{table}

\begin{table}[t]
\caption{Scorer triangulation (full coverage, $n{=}30$). The LLM-judge matches the
unit-tolerant scorer, indicating the headline is scorer-robust.}
\label{tab:scorers}
\begin{tabular}{lc}
\toprule
Scorer & Accuracy \\
\midrule
strict weak-match & 0.20 \\
unit-tolerant & 0.27 \\
LLM-as-judge & 0.27 \\
\bottomrule
\end{tabular}
\end{table}

\textbf{Main finding (Table~\ref{tab:main}).} At fixed model size and identical scorer,
\textsc{GEP} answers 43\% of questions and is 62\% correct on those, versus 20\% for
raw-evidence prompting---a $\sim$3$\times$ improvement in the committed regime, with the
selective curve (Table~\ref{tab:selective}) reaching 0.83 at 20\% coverage.

\textbf{Where it works (Table~\ref{tab:tasktype}).} Gains concentrate on lookup/computation;
committed failures are yes/no and comparison questions whose gold answers are not single scalars
(task-form mismatch), not computation errors---an explicit scope statement, not a hidden failure.

\textbf{Scorer robustness (Table~\ref{tab:scorers}).} Strict, unit-tolerant, and LLM-judge scores
agree on committed items, so the headline is not an artifact of a lenient scorer.

\subsection{Generalization across model families}
\begin{table}[t]
\caption{\textsc{GEP} across model families ($n{=}50$, unit-tolerant). Acc@commit is
accuracy on answered items; coverage is the answered fraction. The effect is not
Qwen-specific: 3 of 4 families reach 0.40--0.56. Nemotron-4B commits often but with wrong
formulas (a calibration failure). Baselines pending except Qwen.}
\label{tab:models}
\begin{tabular}{llccc}
\toprule
Model & Family & Baseline & Acc@commit & Cov. \\
\midrule
Qwen2.5-7B & Qwen & 0.22 & \textbf{0.56} & 0.32 \\
Phi-3.5-mini & Phi & \tbd & \textbf{0.50} & 0.20 \\
deepseek-llm-7b-chat & DeepSeek & \tbd & 0.40 & 0.20 \\
Nemotron-Mini-4B & Nemotron & \tbd & 0.05 & 0.38 \\
\bottomrule
\end{tabular}
\end{table}
Across four families, three reach 0.40--0.56 acc@commit, versus a Qwen raw baseline of 0.22---
evidence the gain is a property of the decomposed procedure, not a single model. Coverage tracks
extractor capability; the weakest extractor (Nemotron-4B) commits frequently to incorrect
programs, motivating the calibrated abstention the method provides. A failure common to all
families is the comparison question type, whose gold answer is not a single scalar.

\subsection{Distillation into a small student (pending)}
\begin{table}[t]
\caption{Held-out reliability of a small student under supervision variants. Metrics:
accuracy, ECE, counterfactual accuracy, teacher-gap recovery.}
\label{tab:distill}
\begin{tabular}{lcccc}
\toprule
Supervision & Acc. & ECE & CF-Acc & Gap-rec. \\
\midrule
answer-only (baseline) & \tbd & \tbd & \tbd & \tbd \\
+ rationale & \tbd & \tbd & \tbd & \tbd \\
+ tool-trace & \tbd & \tbd & \tbd & \tbd \\
\textbf{+ tool-trace + abstain (ours)} & \tbd & \tbd & \tbd & \tbd \\
teacher (reference) & \tbd & \tbd & \tbd & \tbd \\
\bottomrule
\end{tabular}
\end{table}
This is the central method claim: a small student trained to emit tool-verified, calibrated,
abstaining decision programs should exceed answer-only and rationale distillation on accuracy,
calibration, and counterfactual robustness. Results \tbd{} (requires GPU training).

\section{Limitations}\label{sec:limitations}
(i) \textbf{Scope:} the scalar extract--compute procedure targets computational/lookup questions;
selection, list, and open-ended textual questions are out of scope and reported as such.
(ii) \textbf{Generation instability:} fp16 decoding of the 7B extractor on Apple MPS occasionally
degenerates ($\sim$17\% of the pilot), an infrastructure artifact we mitigate but do not remove.
(iii) \textbf{Scale:} pilot numbers are $n{=}30$; full-set ($n{=}150$), multi-seed intervals,
multi-model, and the distillation study are pending (\tbd).

\section{Conclusion}
We reframe small-model financial QA as a decision-extraction-and-computation problem rather than
an access problem, and show that a grounded, tool-verified, selectively-abstaining procedure
yields a $\sim$3$\times$ accuracy gain at fixed model size on computational questions, robust
across scorers. The distillation of this behavior into small students is the natural next step.

\begin{thebibliography}{9}
\bibitem{financebench} P.~Islam et al. FinanceBench: A New Benchmark for Financial Question Answering. 2023.
\bibitem{pot} W.~Chen et al. Program of Thoughts: Disentangling Computation from Reasoning. 2022.
\bibitem{selfconsistency} X.~Wang et al. Self-Consistency Improves Chain-of-Thought Reasoning. 2022.
\bibitem{calibration} K.~Tian et al. Just Ask for Calibration: Eliciting Confidence from LLMs. 2023.
\bibitem{selective} Y.~Geifman and R.~El-Yaniv. Selective Classification for Deep Neural Networks. 2017.
\bibitem{knowledgeconflict} S.~Longpre et al. Entity-Based Knowledge Conflicts in Question Answering. 2021.
\bibitem{distill} C.-Y.~Hsieh et al. Distilling Step-by-Step. 2023.
\end{thebibliography}

\end{document}
